%% 段をまたぐフロートとgyoudoriのfloatオプション（位置を省いたもの）
\documentclass[twocolumn,paper=b5,fontsize=9pt]{jlreq}
\usepackage[float]{gyoudori}\usepackage[columns=3,debug]{evenend}
\setlength\columnseprule{0.4pt}
\newcommand\para{吾輩は猫である．名前はまだ無い．どこで生れたかとんと見当がつかぬ．何でも薄暗いじめじめした所でニャーニャー泣いていた事だけは記憶している．吾輩はここで始めて人間というものを見た．\par}
\newcommand\fig[2]{\centering\fbox{\parbox[c][#2][c]{.95\linewidth}{\centering #1}}\caption{#1}}
\begin{document}
\section{2段にまたがる図}
\para\para
\begin{spanfigure}[b]{2}
\fig{2段にまたがる図（下）}{60pt}
\end{spanfigure}
\para\para\para\para\para\para\para\para\para\para\para\para\para\para\para\para
\begin{spanfigure}{2}
\fig{2段にまたがる図（上）}{50pt}
\end{spanfigure}
\begin{figure}[t]
\fig{1段の図}{40pt}
\end{figure}
\para\para\para\para\para\para\para\para\para\para\para\para\para\para\para\para\para\para\para\para\para\para\para\para
\end{document}
